\section{Drehzahlregelung der Gleichstrommaschine}

Diese Section behandelt die \textbf{Drehzahlregelung der Gleichstrommaschine}
im Zusammenspiel mit dem \textbf{gesteuerten Gleichrichter (B2C)} oder dem \textbf{Gleichstromsteller}.
Ziel ist die systematische Analyse von \textbf{Regulierkennlinien}, \textbf{Betriebspunkten},
\textbf{Grenzbetrieb} und \textbf{Stabilität}.

\smallskip

\textbf{Zielbild:}
\begin{itemize}
    \item Zusammenhang zwischen Zündwinkel / Tastgrad und Drehzahl verstehen
    \item Regulierkennlinien ableiten
    \item Stabilen und instabilen Betrieb erkennen
    \item Grenzfälle des Regelbereichs beurteilen
\end{itemize}

\subsection{Grundstruktur der Antriebsregelung}

Typische Kette:
\[
\text{Netz} \rightarrow \text{Gleichrichter / Steller} \rightarrow \text{Gleichstrommaschine} \rightarrow \text{Last}.
\]

Die Regelgrösse ist:
\begin{itemize}
    \item beim B2C: Zündwinkel $\alpha$  
    \item beim Chopper: Tastgrad $d$
\end{itemize}

Die Stellgrösse ist die \textbf{Ankerspannung}:
\[
\cbl{u_a = \overline{u_2}.}
\]

\subsection{Stationäre Grundgleichung}

Aus der Maschinengleichung:
\[
\cbl{u_a = R_a i_a + \omega \Psi.}
\]

Auflösen nach der Drehzahl:
\[
\cbl{\omega = \frac{u_a - R_a i_a}{\Psi}.}
\]

Interpretation:
\begin{itemize}
    \item Drehzahl steigt mit $u_a$.
    \item Drehzahl sinkt mit steigendem Laststrom.
\end{itemize}

\subsection{Regulierkennlinien}

Für eine feste Lastkennlinie $M_L(\omega)$
und variable Ankerspannung entstehen \textbf{Regulierkennlinien}.

Mit:
\[
M = i_a \Psi, \qquad i_a = \frac{M}{\Psi}
\]

folgt:
\[
\cbl{\omega = \frac{u_a}{\Psi} - \frac{R_a}{\Psi^2} M.}
\]

Das ist eine \textbf{Geradenschar}.



\textbf{Aussage:}
\begin{itemize}
    \item Jede Spannung $u_a$ erzeugt eine eigene Kennlinie.
    \item Regelung erfolgt durch \textbf{Verschieben der Kennlinie}.
\end{itemize}

\subsection{Arbeitspunktbildung}

Der Arbeitspunkt ergibt sich aus:
\[
\omega(M) = \omega_L(M).
\]

Graphisch: Schnittpunkt von
\begin{itemize}
    \item Motorkennlinie
    \item Lastkennlinie
\end{itemize}



\textbf{Aussage:}
\begin{itemize}
    \item Änderung von $u_a$ verschiebt den Arbeitspunkt.
    \item Lastkennlinie bleibt unverändert.
\end{itemize}

\subsection{Einfluss des Zündwinkels $\alpha$ (B2C)}

Beim gesteuerten Gleichrichter gilt:
\[
\cbl{u_a(\alpha) = \frac{2\hat U}{\pi} \cos(\alpha).}
\]

Damit:
\[
\cbl{\omega(\alpha) = \frac{1}{\Psi}
\left( \frac{2\hat U}{\pi} \cos(\alpha) - R_a i_a \right).}
\]

Interpretation:
\begin{itemize}
    \item $\alpha \uparrow \Rightarrow u_a \downarrow \Rightarrow \omega \downarrow$.
    \item $\alpha = \pi/2 \Rightarrow u_a = 0 \Rightarrow Stillstand (idealisert)$.
\end{itemize}

\subsection{Einfluss des Tastgrads $d$ (Chopper)}

Beim Gleichstromsteller gilt:
\[
\cbl{u_a(d) = d \, U_{\mathrm{DC}}.}
\]

Damit:
\[
\cbl{\omega(d) = \frac{d\,U_{\mathrm{DC}} - R_a i_a}{\Psi}.}
\]

\subsection{Grenzbetrieb}

\subsubsection{Leerlauf}

\[
i_a \approx 0 \quad \Rightarrow \quad
\cbl{\omega_0 = \frac{u_a}{\Psi}.}
\]

\subsubsection{Kurzschluss / Blockieren}

\[
\omega = 0 \quad \Rightarrow \quad
\cbl{i_{a,\max} = \frac{u_a}{R_a}, \qquad M_{\max} = \frac{u_a \Psi}{R_a}.}
\]

\subsection{Stabilität der Drehzahlregelung}

Stabilitätskriterium:

\[
\cbl{\frac{d\omega_M}{dM} < \frac{d\omega_L}{dM}
\quad \Rightarrow \quad \text{stabil}.}
\]

Interpretation:
\begin{itemize}
    \item Motorkennlinie flacher als Lastkennlinie.
    \item Kleine Störungen klingen ab.
\end{itemize}

\crd{\textrightarrow\ Bei falscher Steigung entsteht \textbf{selbstverstärkende Abweichung}.}

\subsection{Regulierkennlinienfeld}

Für alle zulässigen $\alpha$ oder $d$ entsteht ein \textbf{Regelbereich}.



\textbf{Aussage:}
\begin{itemize}
    \item Nicht alle Kombinationen $(M,\omega)$ sind erreichbar.
    \item Regelgrenzen durch $\alpha \in [0,\pi]$ oder $d \in [0,1]$.
\end{itemize}

\subsection{Typische Merksätze (Prüfungskern)}

\begin{itemize}
    \item \cbl{$u_a(\alpha) = \dfrac{2\hat U}{\pi} \cos(\alpha)$}
    \item \cbl{$u_a(d) = d\,U_{\mathrm{DC}}$}
    \item \cbl{$\omega = \dfrac{u_a - R_a i_a}{\Psi}$}
    \item Drehzahl wird über Spannung geregelt.
    \item Drehmoment wird über Strom geregelt.
\end{itemize}

\subsection{Typische Fehlerquellen}

\begin{itemize}
    \item Lastkennlinie mit Motorkennlinie verwechselt.
    \item Stabilitätskriterium falsch interpretiert.
    \item Grenzstrom nicht berücksichtigt.
    \item $\alpha$-Regelbereich überschritten.
\end{itemize}

\subsection{Minimal-Checkliste für Aufgaben}

\begin{enumerate}
    \item Art des Umrichters bestimmen (B2C oder Chopper).
    \item Stellgrösse identifizieren ($\alpha$ oder $d$).
    \item $u_a$ berechnen.
    \item Kennlinie $\omega(M)$ aufstellen.
    \item Arbeitspunkt mit Last bestimmen.
    \item Stabilität qualitativ prüfen.
\end{enumerate}
